% !TEX root = ../my-thesis.tex
%
\chapter{Evaluation}
\label{sec:evaluation}

This chapter evaluates the pipeline. Section~\ref{sec:evaluation:setup} describes
the experimental setup. Section~\ref{sec:evaluation:crosstopic} reports the
cross-topic results over the eight UKP topics (RQ1, RQ4), and
Section~\ref{sec:evaluation:llm} adds a qualitative check on LLM-generated
arguments. Section~\ref{sec:evaluation:ablation} compares typed against binary
prompting (RQ2), and Section~\ref{sec:evaluation:gold} describes the
evaluation of the attack quality against a manual reference annotation (RQ3).
Section~\ref{sec:evaluation:case} closes the chapter with a small case study
and an analysis of the recurring error types.

\section{Experimental Setup}
\label{sec:evaluation:setup}

The pipeline is evaluated in two settings. The first starts from argument sets of
the UKP Sentential Argument Mining Corpus~\cite{stab2018}, which provides
sentence-level stance annotations for eight controversial topics. Each sentence
is labelled \texttt{Argument\_for}, \texttt{Argument\_against}, or
\texttt{NoArgument}. The \texttt{NoArgument} sentences are discarded and the
test splits are used. Table~\ref{tab:evaluation:corpus} lists the material
available per topic. From each topic a balanced sample of 14 arguments, seven
pro and seven con, is drawn with the fixed random seed 42, so that the sample
is reproducible. The second setting covers the case the pipeline ultimately
targets, arguments produced by an LLM, and uses LLM-generated arguments on the
topic of allowing autonomous vehicles in public, with ten pro and ten con
arguments. Since the generation prompt already constrains each argument to a
single claim, atomization is skipped in that setting. The corpus setting
carries the quantitative evaluation, because its arguments are not produced by
the model under test and therefore provide an independent reference, while the
generated setting is a qualitative demonstration that the same pipeline runs
on LLM-produced input, as Section~\ref{sec:evaluation:llm} discusses.

\begin{table}[tb]
  \centering
  \caption{Arguments available in the UKP test splits after discarding
  \texttt{NoArgument} sentences. The evaluation samples seven pro and seven con
  arguments per topic with seed 42.}
  \label{tab:evaluation:corpus}
  % Auto-generiert von code/make_thesis_tables.py — nicht von Hand editieren.
\begin{tabular}{lrrr}
  \toprule
  Topic & Pro & Con & Total \\
  \midrule
  Abortion          & 136 & 165 & 301 \\
  Cloning           & 142 & 168 & 310 \\
  Death Penalty     & 103 & 232 & 335 \\
  Gun Control       & 158 & 133 & 291 \\
  Marijuana Legal.  & 118 & 126 & 244 \\
  Minimum Wage      & 116 & 111 & 227 \\
  Nuclear Energy    & 122 & 171 & 293 \\
  School Uniforms   & 109 & 146 & 255 \\
  \bottomrule
\end{tabular}

\end{table}

In both settings the extractor and the verifier call
GPT-OSS-120B~\cite{gptoss2025}, an open-weight mixture-of-experts model with
117 billion parameters, through a university-hosted endpoint at temperature
zero, with argument pairs processed in batches of $b = 10$. Candidates are accepted at
confidence $c \geq \theta_{\text{acc}} = 70$, forwarded to the verifier at
$\theta_{\text{ver}} \leq c < \theta_{\text{acc}}$ with
$\theta_{\text{ver}} = 50$, and rejected below $\theta_{\text{ver}}$. Since
the model is open weight, the experiments do not depend on a proprietary
service remaining available or unchanged, and the same weights can in
principle be redeployed by anyone.

The reproducibility discussion in Section~\ref{sec:foundations:llm} shapes
the setup in two ways. Atomization is not perfectly deterministic even at
temperature zero and produced 21 atoms in one run and 22 in another for the
same topic despite identical inputs and seed, so each topic is atomized
exactly once and all experiments on it, including both arms of the ablation
and the reference annotation, operate on the same frozen atom set. For the same
reason, the cross-topic and ablation numbers stem from one canonical batch of
runs with one prompt version, and the reference evaluation scores of
Section~\ref{sec:evaluation:gold} from one dedicated run with the final
prompt revision, so that no table mixes results from different prompt
wordings. The logging infrastructure of Section~\ref{sec:method:logging}
ties every reported number to a concrete run.

\section{Cross-Topic Results}
\label{sec:evaluation:crosstopic}

Table~\ref{tab:evaluation:crosstopic} reports, for each of the eight UKP topics,
the number of atomic arguments after atomization, the number of queried pairs,
the number of accepted attacks and their share among all queried pairs, the
number of isolated atoms without any incident attack, the size of the grounded
extension, the size of the undecided core $|U|$ from
Section~\ref{sec:method:solver}, and the numbers of preferred and stable
extensions.

\begin{table}[tb]
  \centering
  \caption{Cross-topic results over the eight UKP topics. ``Atoms'' is the
  number of atomic arguments after atomization, and ``Pairs'' is the number of
  queried pairs, excluding pairs of atoms that stem from the same source
  argument. ``Acc'' is the number
  of accepted attacks, that is the edges of the framework, and ``Acc\%'' is
  their share among all queried pairs. ``Iso'' counts atoms without incident
  attacks, ``Grnd'' is the size of the grounded extension, $|U|$ the size of
  the undecided core, and \#Pr and \#St the numbers of preferred and stable
  extensions.}
  \label{tab:evaluation:crosstopic}
  \small
  % Auto-generiert von code/make_thesis_tables.py — nicht von Hand editieren.
\begin{tabular}{lrrrrrrrrr}
  \toprule
  Topic & Atoms & Pairs & Acc & Acc\% & Iso & Grnd & $|U|$ & \#Pr & \#St \\
  \midrule
  Abortion          & 25 & 267 & 41 & 15.4 & 6 & 7 & 9 & 4 & 4 \\
  Cloning           & 20 & 180 & 59 & 32.8 & 2 & 9 & 3 & 2 & 2 \\
  Death Penalty     & 18 & 148 & 68 & 45.9 & 1 & 3 & 8 & 3 & 3 \\
  Gun Control       & 21 & 203 & 26 & 12.8 & 5 & 14 & 0 & 1 & 1 \\
  Marijuana Legal.  & 20 & 182 & 41 & 22.5 & 2 & 13 & 0 & 1 & 1 \\
  Minimum Wage      & 22 & 221 & 58 & 26.2 & 3 & 10 & 3 & 2 & 2 \\
  Nuclear Energy    & 21 & 201 & 43 & 21.4 & 1 & 12 & 3 & 2 & 2 \\
  School Uniforms   & 21 & 201 & 29 & 14.4 & 3 & 13 & 0 & 1 & 1 \\
  \bottomrule
\end{tabular}

\end{table}

Three observations stand out. Every framework contains cycles, a direct
consequence of the rebuttal symmetrization described in
Section~\ref{sec:foundations:typology}, since every topic has at least three
symmetrized rebuttal edges (column ``Sym'' of
Table~\ref{tab:evaluation:types}) and each symmetrized pair forms a two-cycle.
Five of the eight topics admit
several preferred extensions, and the number of stable extensions matches the
number of preferred extensions in every case. So although
Section~\ref{sec:foundations:aaf} shows that stable extensions can fail to
exist, here every preferred extension is also stable. Each maximal position
attacks everything it does not contain, and none of the extracted frameworks
behaves like the odd-length cycle of Figure~\ref{fig:foundations:example}, whose
undecided conflicts leave no stable extension. These multiple positions are not
asserted by a free generation. They are derived from the model's own pairwise
attack judgments through the fixed semantics, and each is conflict-free and
admissible by construction. The undecided core, finally, stays small
everywhere, with at most nine arguments, which confirms the sparsity
assumption behind the grounded-first reduction. On three topics the grounded extension already
decides every argument and $U$ is empty, so the formal layer adds nothing
beyond the cautious core there. The acceptance rate varies widely, from 12.8
percent for Gun Control to 45.9 percent for Death Penalty, which reflects
genuine differences in how directly the sampled arguments engage each other.

A closer look at Abortion, the topic with the most extensions, shows what the
multiple positions look like. All four preferred extensions share the grounded
core of seven arguments and differ only on the undecided core of nine. Notably,
the extensions do not split cleanly along the pro and con labels of the corpus.
One extension combines several con arguments about ending abortion through
legislation with a pro argument that concedes state regulation after fetal
viability, while another accepts only additional pro arguments. The formal
layer thus recovers positions that are finer-grained than the binary stance
annotation, separating, for instance, a legally framed camp from a
morally framed one. At the same time, two of the undecided arguments are
defective atoms in the sense of Section~\ref{sec:evaluation:errors}, so part of
the multiplicity on this topic rests on atoms that a cleaner atomization would
have removed. Both faces of the result, genuinely finer positions and
inflation through defective atoms, recur throughout the error analysis.

\begin{table}[tb]
  \centering
  \caption{Composition of the accepted attacks per topic. ``Reb'', ``Und'', and
  ``Cut'' count edges labelled rebuttal, undermining, and undercutting,
  including symmetrized reverse edges. ``Sym'' is the number of edges added by
  rebuttal symmetrization and is contained in ``Reb''. ``Ver'' counts
  candidates that were accepted only after verification. ``Extract'' is the
  wall-clock extraction time in seconds. The three type counts per topic sum
  to the Acc column of Table~\ref{tab:evaluation:crosstopic}.}
  \label{tab:evaluation:types}
  % Auto-generiert von code/make_thesis_tables.py — nicht von Hand editieren.
\begin{tabular}{lrrrrrr}
  \toprule
  Topic & Reb & Und & Cut & Sym & Ver & Extract (s) \\
  \midrule
  Abortion          & 30 & 2 & 9 & 15 & 0 & 337 \\
  Cloning           & 14 & 35 & 10 & 7 & 0 & 161 \\
  Death Penalty     & 22 & 41 & 5 & 11 & 1 & 233 \\
  Gun Control       & 6 & 19 & 1 & 3 & 0 & 335 \\
  Marijuana Legal.  & 20 & 6 & 15 & 10 & 1 & 243 \\
  Minimum Wage      & 6 & 37 & 15 & 3 & 0 & 434 \\
  Nuclear Energy    & 12 & 29 & 2 & 6 & 1 & 148 \\
  School Uniforms   & 10 & 17 & 2 & 5 & 3 & 173 \\
  \bottomrule
\end{tabular}

\end{table}

Table~\ref{tab:evaluation:types} breaks the accepted attacks down by type.
Undermining is the most common label overall, which fits the corpus, since many
UKP sentences report evidence and statistics that target the premises of
opposing claims rather than negating their conclusions outright. Undercutting
is the rarest label on most topics. The distribution varies strongly by topic.
Abortion and Marijuana Legalization are rebuttal-heavy, and on Abortion 15 of
the 41 edges exist only through symmetrization, which foreshadows the error
analysis in Section~\ref{sec:evaluation:errors}. Verification adds few edges.
Across all eight topics, 68 of the 1603 queried pairs were routed to the
verifier, which rejected 62 of them and confirmed six, so the extractor's
confidence scores separate clear accepts from
clear rejects for about 96 percent of the queried pairs. Extraction takes between
two and eight minutes per topic and is dominated by API latency, while solving
all three semantics takes under half a second per topic thanks to the
grounded-first reduction.

\section{The LLM-Generated Setting}
\label{sec:evaluation:llm}

The second setting feeds twenty LLM-generated arguments on autonomous vehicles
through the same extraction and solving stages. Of the 190 queried pairs, 37
are accepted as attacks, an acceptance rate of 19.5 percent that lies within
the range of the UKP topics in Table~\ref{tab:evaluation:crosstopic}. The
resulting framework is cyclic, yet the grounded extension already accepts or
attacks every argument, so all three semantics coincide on a single extension
of ten arguments. It contains eight of the ten con arguments but only the two
pro arguments that participate in no attack at all, so the extraction judges
the con side dominant on this generated argument set. The setting is a
qualitative check rather than a second evaluation, since the arguments, the
attacks, and the confidence scores all originate from the same model and the
run predates the final prompt revision. It shows that the pipeline runs
end-to-end on generated input and that generated argument sets, which are cleaner and
more uniform than corpus sentences, tend to be decided completely by the
grounded extension.

\section{Typed versus Binary Prompting}
\label{sec:evaluation:ablation}

This section measures whether the typed attack prompting of
Section~\ref{sec:method:extraction} actually produces different attack
relations than plain binary prompting, and what the difference does to the
resulting frameworks. A second extraction arm runs the identical frozen atom
set $\mathcal{A}'$ of every topic through a binary prompt that asks only
whether an attack exists and in which direction, without the attack type,
without the type definitions of Section~\ref{sec:foundations:typology}, and
without the two rebuttal rules of Section~\ref{sec:method:extraction}. Since the binary arm produces no type
labels, no rebuttal symmetrization is possible there. To separate the effect of
the categorization from the effect of symmetrization, a third arm is evaluated
that takes the typed attack set but drops the symmetrized reverse edges. This
arm, called typed-raw, requires no additional extraction, only a second solver
run on the reduced edge set.

The overlap of two edge sets $E_1$ and $E_2$ is measured by their
\emph{Jaccard similarity}
\begin{equation*}
J(E_1, E_2) = \frac{|E_1 \cap E_2|}{|E_1 \cup E_2|},
\end{equation*}
which is 1 for identical and 0 for disjoint sets. It is computed twice, once
on the directed edges and once on undirected conflict pairs, where each edge
is collapsed to the unordered pair of its endpoints. The undirected variant
asks the weaker question whether the two arms mark the same pairs as
conflicting at all, regardless of the asserted direction.

\begin{table}[tb]
  \centering
  \caption{Typed versus binary prompting on identical atom sets. ``$|E|$''
  columns give the number of accepted directed edges for the typed arm, the
  typed arm without symmetrized edges (raw), and the binary arm.
  ``$J$'' columns give the Jaccard similarity of typed-raw and binary edge
  sets, directed and undirected. \#Pr counts preferred extensions per arm, and
  the numbers of stable extensions are identical to \#Pr in every cell.
  ``Same-st.'' counts accepted attacks between arguments of the same stance,
  for the typed arm on the final edge set including symmetrized edges.}
  \label{tab:evaluation:ablation}
  \small
  \setlength{\tabcolsep}{4pt}
  % Auto-generiert von code/make_thesis_tables.py — nicht von Hand editieren.
\begin{tabular}{lrrrrrrrrrr}
  \toprule
  & \multicolumn{3}{c}{$|E|$} & \multicolumn{2}{c}{$J$}
  & \multicolumn{3}{c}{\#Pr} & \multicolumn{2}{c}{Same-st.} \\
  \cmidrule(lr){2-4}\cmidrule(lr){5-6}\cmidrule(lr){7-9}\cmidrule(lr){10-11}
  Topic & typ. & raw & bin. & dir. & undir. & typ. & raw & bin. & typ. & bin. \\
  \midrule
  Abortion          & 41 & 26 & 39 & 0.08 & 0.12 & 4 & 1 & 1 & 3 & 14 \\
  Cloning           & 59 & 52 & 35 & 0.47 & 0.47 & 2 & 1 & 1 & 7 & 1 \\
  Death Penalty     & 68 & 57 & 51 & 0.71 & 0.83 & 3 & 1 & 1 & 1 & 2 \\
  Gun Control       & 26 & 23 & 23 & 0.35 & 0.39 & 1 & 1 & 1 & 1 & 1 \\
  Marijuana Legal.  & 41 & 31 & 24 & 0.41 & 0.53 & 1 & 1 & 1 & 1 & 0 \\
  Minimum Wage      & 58 & 55 & 49 & 0.27 & 0.44 & 2 & 1 & 1 & 0 & 10 \\
  Nuclear Energy    & 43 & 37 & 27 & 0.56 & 0.64 & 2 & 1 & 1 & 0 & 0 \\
  School Uniforms   & 29 & 24 & 42 & 0.38 & 0.38 & 1 & 1 & 1 & 0 & 1 \\
  \bottomrule
\end{tabular}

\end{table}

Table~\ref{tab:evaluation:ablation} summarises the comparison. Three findings
emerge.

\paragraph{The two prompts extract substantially different attack sets.}
The directed Jaccard similarity between the typed-raw and binary edge sets
ranges from 0.08 to 0.71 with a mean around 0.4. Even on Death Penalty, the
topic with the highest agreement, more than a quarter of the union of both edge
sets is found by only one arm. On Abortion the two arms are nearly disjoint,
with only five shared edges out of 26 typed-raw and 39 binary edges. The typed
arm accepts more raw edges on five of the eight topics, the two arms tie on
Gun Control, and the binary arm
accepts more on Abortion and School Uniforms. Asking for the attack type
therefore changes which conflicts the model asserts rather than rephrasing the
same judgment.

\paragraph{Multiple extensions come from symmetrization, not from typing
itself.}
On the final typed frameworks, five of eight topics admit several preferred
extensions. On the typed-raw frameworks, every topic collapses to exactly one
preferred and one stable extension, the same count as the binary arm. The
multiplicity of positions that Section~\ref{sec:evaluation:crosstopic} reports
is thus produced entirely by the mutual attack edges that rebuttal
symmetrization adds, a mechanism that is only available when the extraction
knows which edges are rebuttals. What this dependence on one modelling
decision means for the value of the formal layer is discussed in
Section~\ref{sec:discussion:value}.

\paragraph{The arms differ in their characteristic false positives.}
Without a manual reference annotation, edge counts alone cannot say which arm judges more
accurately, but the accepted same-stance attacks give a diagnostic signal.
Attacks between arguments of the same stance are possible in principle, yet
they should be rare in balanced debate samples. The binary arm accepts 14
same-stance attacks on Abortion and 10 on Minimum Wage, among them a cluster in
which two statistical claims about abortion frequency attack seven separate
contra claims about post-abortion experiences, an implausible pattern within
one stance camp. The typed arm accepts far fewer same-stance attacks on those
topics, which suggests that forcing the model to name the attack type filters
out part of this over-triggering. The signal is not uniform, though. On
Cloning the relation reverses, with seven typed against one binary same-stance
attack, so the same-stance count is a topic-dependent indicator rather than a
general proof of the typed advantage. The typed arm has its own failure mode, as
Section~\ref{sec:evaluation:errors} shows, since several of its Abortion edges
attack atoms that are not genuine arguments in the first place. Which of the
two error profiles weighs heavier is exactly the question the manual annotation of
the next section is designed to answer.

\section{Attack Quality against a Manual Reference Annotation}
\label{sec:evaluation:gold}

The ablation shows that typed and binary prompting extract different attack
sets, but it cannot decide which set is closer to the truth. To answer that, a
manually annotated reference set is constructed on the Nuclear Energy topic. The frozen
atom set of the canonical run, 21 atoms, yields 201 argument pairs after
excluding pairs that share a source argument, exactly the pair set the
extraction judges. Every pair is annotated by hand with one of four labels,
namely no attack, A attacks B, B attacks A, or mutual attack, following written
guidelines based on the attack definitions of Pollock and
ASPIC+~\cite{pollock1992, modgil2014}. A validation script checks that the
annotated pair set matches the queried pair set exactly and that all labels are
well-formed. The annotation records only the existence and direction of an
attack, not its type. Assigning types by hand would have required
premise-level judgments for every pair from a single annotator, and the
reference is meant to measure attack detection, on which the typed and the
binary arm can be compared directly, rather than type assignment.

The guidelines instruct the annotator to settle the attack question before the
direction question and to judge every claim for the object it names as
written. A safety claim about one subsystem, for example waste storage, is not
read as a claim about the whole nuclear operation, so a conflict that only
arises after widening a claim to an object it does not name is labelled as no
attack. An attack may rest on a minimal implicit premise, but not on such an
object switch. Borderline decisions are recorded in a note column and revisited
for consistency before the labels are frozen. The final annotation labels 190
pairs as no attack, 6 pairs as one-directional attacks, and 5 pairs as mutual
attacks, which yields 16 directed reference edges and 11 undirected conflict pairs.
Genuine conflicts are therefore rare, only about one pair in eighteen carries
an attack at all, and this sparsity is worth keeping in mind when reading the
scores, since every single reference edge moves recall by six points.

Both extraction arms are then scored against the reference labels, so that the
reference annotation simultaneously provides the quality measurement for RQ3 and the
missing reliability comparison for RQ2. An extracted attack that the reference
annotation confirms is a true positive (TP), an extracted attack that the reference
annotation rejects is a false positive (FP), and an annotated attack that the
extraction misses is a false negative (FN). The evaluation reports

\begin{equation*}
\text{precision} = \frac{TP}{TP + FP}, \qquad
\text{recall} = \frac{TP}{TP + FN}, \qquad
\text{F1} = \frac{2 \cdot \text{precision} \cdot \text{recall}}
                 {\text{precision} + \text{recall}},
\end{equation*}

\noindent
that is, the share of extracted attacks that are genuine, the share of genuine
attacks that are found, and their harmonic mean. All three are computed on the
attack class, both for directed edges and for undirected conflict pairs.
Following the metric argument of Section~\ref{sec:related:ukp}, accuracy is
deliberately not reported. About 46 percent of the 201 pairs connect arguments
of the same stance, and since genuine same-stance attacks are rare, a majority
of easy negatives would inflate accuracy without saying anything about the
quality of the extracted attacks. Precision, recall, and F1 do not contain the
true negatives in their formulas, so including the same-stance pairs makes the
evaluation stricter rather than more lenient, because every accepted
same-stance attack that the reference annotation rejects counts as a false positive.

The scored predictions come from a dedicated evaluation run that executes both
extraction arms on the identical frozen atom set with the final revision of
the typed and verifier prompts. This revision, made after the ablation batch
of Section~\ref{sec:evaluation:ablation}, adds an explicit test that two
claims which can both be true without tension do not attack each other,
whatever their stances, and it requires an undermining attack to name the
premise it hits. The typed arm accepts 23
directed edges in this run against 43 in the ablation batch, which is why the
edge counts differ from the Nuclear Energy row of
Table~\ref{tab:evaluation:ablation}. The binary prompt is identical in both
batches, and its count still moves from 27 to 32 edges through the
run-to-run nondeterminism discussed in Section~\ref{sec:evaluation:setup}.
The first difference illustrates the prompt sensitivity and the second the
residual nondeterminism that Section~\ref{sec:discussion:threats} discusses
among the limitations. Scoring happens on
three levels. The directed final level compares all accepted edges including
the symmetrized reverse edges against the 16 directed reference edges. The directed
raw level drops the symmetrized edges and asks how well the model asserts
directions on its own. The undirected level collapses edges to conflict pairs
and compares them against the 11 reference pairs. Since the binary arm produces no
type labels, its raw and final edge sets coincide.
Table~\ref{tab:evaluation:gold} shows the results.

\begin{table}[tb]
  \centering
  \caption{Precision, recall, and F1 of both extraction arms against the
  manual reference annotation on Nuclear Energy. ``$n$'' is the number of predicted edges or
  pairs, TP the number of true positives. The reference annotation contains 16
  directed edges and 11 undirected conflict pairs.}
  \label{tab:evaluation:gold}
  \small
  % Auto-generiert von code/make_thesis_tables.py — nicht von Hand editieren.
\begin{tabular}{llrrrrr}
  \toprule
  Arm & Level & $n$ & TP & Precision & Recall & F1 \\
  \midrule
  typed  & directed final    & 23 & 12 & 0.522 & 0.750 & 0.615 \\
         & directed raw      & 19 & 9 & 0.474 & 0.562 & 0.514 \\
         & undirected pairs  & 19 & 9 & 0.474 & 0.818 & 0.600 \\
  \midrule
  binary & directed          & 32 & 11 & 0.344 & 0.688 & 0.458 \\
         & undirected pairs  & 32 & 11 & 0.344 & 1.000 & 0.512 \\
  \bottomrule
\end{tabular}

\end{table}

\paragraph{Overall performance.}
The reference evaluation resolves the question that the ablation left open. The typed
arm outperforms the binary arm on every level, with a directed F1 of 0.615
against 0.458. The difference is driven by precision. The binary arm finds
every annotated conflict pair, a recall of 1.000 on the undirected level, but it
asserts almost three times as many conflict pairs as the reference annotation
contains, so the 11 genuine conflicts are buried among 21 spurious ones. The
typed arm asserts fewer edges and a larger share of them is genuine. This
confirms by direct measurement what the same-stance diagnostic of
Section~\ref{sec:evaluation:ablation} could only suggest, namely that requiring
the model to name an attack type acts as a precision filter. The symmetrized
rebuttal edges also earn their place, since they lift the typed directed F1
from 0.514 to 0.615 by supplying reverse directions of mutual conflicts that
the model does not assert on its own. The binary arm has no access to this
mechanism, and all five of its missed directed edges are reverse directions
of mutual reference pairs whose forward direction it does assert.

\paragraph{False positives.}
The false positives of both arms follow one dominant pattern. The model
accepts attacks between claims that belong to different sub-topics of the
debate. A
representative case is the accepted undermining attack from the atom
``building a plant with 100\,\% security is technically impossible'' on the
atom ``storage is safe and secure''. The first claim concerns reactor
operation and the second waste storage, so under the strict object reading of
the guidelines the pair is compatible, and an attack only arises after
widening the security claim to the whole nuclear operation. The same widening
connects the water-consumption claims to the claim that nuclear is very
competitive once external costs are accounted. These false positives carry
confidences up to 90, so raising the acceptance threshold would not remove
them without also removing genuine attacks. The verification stage, in
contrast, works in the intended direction, since both candidates it rejected
in the typed run are pairs that the reference annotation also rejects.

\paragraph{Missed attacks.}
All four missed directed edges of the typed arm belong to the two mutual reference
pairs built on the atom ``alternative methods would yield much more
climate-saving bang for our buck than nuclear power, according to some
analysts''. The reference annotation treats this claim as directly contradicting
the atoms ``nuclear is much cheaper than wind or photovoltaic systems'' and
``if external costs are accounted, nuclear is very competitive'', but the
model judges both pairs as compatible in every run. A plausible explanation is
the hedged phrasing, which attributes the comparison to unnamed analysts, so
the model reads the atom as a report rather than an asserted conclusion.
Detecting rebuttals between hedged comparative claims is thus the clearest
remaining recall gap.

\paragraph{Stability and context.}
Three repeated runs of the evaluated configuration on the frozen atom set yield
typed directed F1 scores between 0.579 and 0.615, so the typed advantage over
the binary arm, sixteen points on the evaluated run, clearly exceeds the run to
run variation. For rough context, the scores lie in a similar range as the
sentence-level F1 of 67 percent that the knowledge-based system discussed in
Section~\ref{sec:related:ukp} reaches on the same corpus, on a task that is
related but not directly comparable~\cite{kramer2021}. In answer to RQ3, the pipeline
recovers three quarters of the manually annotated attacks with moderate
precision, its errors are systematic rather than random, and the typed
extraction arm is measurably closer to the manual annotation than the binary
baseline.

\section{Case Study and Error Analysis}
\label{sec:evaluation:case}
\label{sec:evaluation:errors}

To make the pipeline's behaviour and its typical failures concrete, this
section first works through a deliberately small nuclear-energy sample in full
detail and then collects the recurring error types that the case study, the
ablation, and the reference evaluation surface. The sample is a separate illustrative run
with three pro and three con arguments and is not part of the main
configuration of Section~\ref{sec:evaluation:setup}, whose frameworks are too
large to display as a readable graph. Starting from six balanced UKP arguments,
atomization yields eight atomic arguments, three pro and five con. Only the
sentence about meltdown statistics is split, into three separate claims, the
case already shown in Section~\ref{sec:method:atomize}. Of the
$\binom{8}{2} = 28$ pairs, 25 are queried after skipping pairs of atoms that
share a source argument. Six attacks are accepted. Rebuttal symmetrization adds
two reverse edges, so the framework has eight directed attacks, four
undermining and four rebuttals.

The graph centers on the argument that nuclear power merits support because it
is virtually free of CO\textsubscript{2} emissions, which the con side attacks
from several directions. The symmetrized rebuttals make the framework cyclic, so several preferred
and stable sets might be expected. In this case, however, the
grounded extension already terminates by accepting all five con arguments and
rejecting the entire pro side, so grounded, preferred, and stable coincide on a
single extension.

\begin{figure}[tb]
  \centering
  \begin{minipage}[c]{0.50\linewidth}
    \centering
    \begin{tikzpicture}[
      >={Stealth[length=1.8mm]},
      arg/.style={draw, circle, minimum size=6.5mm, inner sep=0.2pt, font=\tiny\ttfamily, fill=white},
      pro/.style={arg, draw=blue!55!black},
      con/.style={arg, draw=red!55!black},
      gr/.style={fill=gray!22},
      reb/.style={red!75!black, <->, thin},
      und/.style={blue!60!black, ->, thin},
    ]
      \node[con,gr] (con58)  at (-1.1, 1.2) {con58};
      \node[con,gr] (con71b) at (-1.6, 0.0) {con71b};
      \node[con,gr] (con71c) at (-1.1,-1.2) {con71c};
      \node[pro]    (pro4)   at ( 0.0, 0.0) {pro4};
      \node[con,gr] (con63)  at ( 2.4, 0.0) {con63};
      \node[pro]    (pro82)  at ( 3.4, 1.0) {pro82};
      \node[pro]    (pro15)  at ( 3.4,-1.0) {pro15};
      \node[con,gr] (con71)  at ( 1.2,-1.9) {con71};
      \draw[reb] (con71b) to (pro4);
      \draw[reb] (con71c) to (pro4);
      \draw[und] (con58) to (pro4);
      \draw[und] (con63) to (pro4);
      \draw[und] (con63) to (pro82);
      \draw[und] (con63) to (pro15);
    \end{tikzpicture}
  \end{minipage}\hfill
  \begin{minipage}[c]{0.48\linewidth}
    \scriptsize
    \begin{tabular}{@{}l@{\ }p{0.76\linewidth}@{}}
      \texttt{pro4:}   & nuclear is $\text{CO}_2$ emission-free, merits support \\
      \texttt{pro82:}  & waste storage is safe and secure \\
      \texttt{pro15:}  & spent-fuel storage is robust \\
      \texttt{con63:}  & nuclear waste is a terrorism risk \\
      \texttt{con58:}  & subsidies better spent on renewables \\
      \texttt{con71b:} & 1 of 5 meltdowns: 1000+ deaths \\
      \texttt{con71c:} & 1 of 100k meltdowns: 50000 deaths \\
      \texttt{con71:}  & 2 of 3 meltdowns: no deaths (isolated) \\
    \end{tabular}
  \end{minipage}
  \caption{Attack graph of the small nuclear-energy sample, eight atoms and eight
  directed edges. Red double arrows are symmetric rebuttals, blue arrows denote
  undermining attacks. Filled nodes form the grounded extension, which here
  coincides with the unique preferred and stable extension.}
  \label{fig:evaluation:nuclear}
\end{figure}

Beyond this small sample, three recurring error types stand out across the
experiments, each originating in a different stage of the pipeline.

\paragraph{Rebuttal mislabeling creates unjustified cycles.}
The attack-type classification can mislabel a premise-directed attack as a
rebuttal. Since rebuttals are symmetrized into mutual attacks while undermining
attacks stay directed, a false rebuttal label introduces a bidirectional
conflict and therefore an unjustified cycle. In the case study, the
meltdown-fatality claims attack the CO\textsubscript{2} argument as rebuttals,
although they question its safety premise rather than directly negating its
support conclusion. An undermining label would have been more suitable and
would not have added a reverse edge. The reference evaluation run of
Section~\ref{sec:evaluation:gold} shows both faces of this mechanism on the
same topic. Three of the four symmetrized edges in that run are confirmed by
the reference annotation, while the fourth turns a one-directional attack among the
cost claims into a spurious mutual conflict and enters
Table~\ref{tab:evaluation:gold} as the false positive with the highest
confidence of the run. The cost of this error type scales with
the number of symmetrized edges, which
Table~\ref{tab:evaluation:types} quantifies. Between 3 and 15 edges per topic
exist only through symmetrization, and on the rebuttal-heavy Abortion topic
more than a third of all edges are symmetrized reverse edges. Every unjustified
mutual attack potentially splits the framework into additional extensions, so
the multiplicity results of Section~\ref{sec:evaluation:crosstopic} must be
read with this caveat in mind. The descriptive-versus-normative rule described
in Section~\ref{sec:method:extraction} was added precisely to curb this error
type, and its effect on the numbers is part of the prompt-sensitivity
discussion in Section~\ref{sec:discussion:threats}.

\paragraph{Atomization produces defective atoms.}
Atomization can produce atoms that, read in isolation, are not self-contained
arguments. Three variants occur. Some atoms lose their stance, such as the
claim that two out of three meltdowns would cause no deaths, which understates
rather than opposes the pro side and stays isolated in the graph. Some atoms keep an
unresolved reference despite the self-containment rule, such as the Abortion
atom ``not all abortions are unjustified according to this argument''. Some
atoms are not argumentative at all, such as a sentence that merely cites the
introduction of a legislative act. Table~\ref{tab:evaluation:crosstopic} gives
an upper bound on the harmless variant, since between 1 and 6 atoms per topic
end up isolated. The harmful variant is when a defective atom does participate
in attacks. On Abortion, several edges found only by the typed arm attack
exactly the two defective atoms above, so extraction errors and atomization
errors compound.

\paragraph{Binary prompting over-triggers within one stance camp.}
The third error type is specific to the binary arm. Without the type
definitions, the model accepts attacks between same-stance arguments far more
liberally, most visibly in the Abortion cluster already described in
Section~\ref{sec:evaluation:ablation}, a pattern that the typed prompt largely
suppresses on the affected topics. This asymmetry of error profiles, spurious
cycles and defective targets on the typed side against same-stance
over-triggering on the binary side, is why edge counts alone could not settle
the comparison, and why the verdict of Section~\ref{sec:evaluation:gold} rests
on the manual annotation rather than on a metric chosen in either arm's favour.
